% GENERATED FILE. Edit canonical lessons, then run npm run generate:books.
% canonical-source-hash: fnv1a64:38ecb778fcef3101
% canonical-lessons: UR-C504-svich, UR-C504-plag, UR-C504-karant, UR-C504-khunti, UR-C504-saket

\chapter{Wires and Switches}
\label{ch:ur-wires-and-switches}
\hlchaptermodality{504}

\begin{chapteropening}
\textbf{By the end of this chapter:} \emph{I can talk about lights, plugs and sockets at home.}

Complete the last lesson of chapter 504: I can talk about lights, plugs and sockets at home.
\end{chapteropening}

\section[\texorpdfstring{\ur{سوئچ} (svich)}{سوئچ (svich)}]{\ur{سوئچ} (svich) — a switch}

\label{lesson:UR-C504-svich}



\begin{quote}
Before the new one: say the Urdu for to have (something) painted, then the Urdu for to smear (with mud).
\end{quote}

\subsection*{You'll want to know: \ur{سوئچ}}
\textbf{\ur{سوئچ}} — \emph{svich} — \textquotedblleft{}a switch\textquotedblright{}.

\begin{quote}
\ur{بتی} \ur{کا} \ur{سوئچ} \ur{دروازے} \ur{کے} \ur{پاس} \ur{ہے۔}
\end{quote}

\begin{quote}
\emph{battī kā svich darvāze ke pās hai.}
\end{quote}

\begin{quote}
The light switch is by the door.
\end{quote}

Masculine.

\begin{cousinweb}
\textbf{\ur{سوئچ}} (\emph{svich}) is borrowed from English \emph{switch}.
\end{cousinweb}

\begin{grammarlens}[title={What you've built}]
One more word for describing what you lived and felt.
\end{grammarlens}

\subsection*{Your turn}
\begin{itemize}
\raggedright
  \item \emph{Say it:} \emph{svich}
  \item \emph{Say it:} \emph{svich}, once more
  \item \emph{Say it:} \emph{battī kā svich darvāze ke pās hai.} — The light switch is by the door.
  \item \emph{Recall:} say the Urdu for to have (something) painted, then the Urdu for to smear (with mud), then say \emph{svich} again
\end{itemize}

\begin{tcolorbox}[breakable,colback=teal!4,colframe=teal!35!black,title={Before you move on}]
What does \ur{سوئچ} mean? (\textquotedblleft{}A switch\textquotedblright{}.) Say it once more.
\end{tcolorbox}
\section[\texorpdfstring{\ur{پلگ} (plag)}{پلگ (plag)}]{\ur{پلگ} (plag) — a plug (on a cable)}

\label{lesson:UR-C504-plag}



\begin{quote}
Before the new one: say the Urdu for to smear (with mud), then the Urdu for a switch.
\end{quote}

\subsection*{You'll want to know: \ur{پلگ}}
\textbf{\ur{پلگ}} — \emph{plag} — \textquotedblleft{}a plug (on a cable)\textquotedblright{}.

\begin{quote}
\ur{اس} \ur{چارجر} \ur{کا} \ur{پلگ} \ur{کہاں} \ur{ہے؟}
\end{quote}

\begin{quote}
\emph{is chārjar kā plag kahāṅ hai?}
\end{quote}

\begin{quote}
Where is the plug for this charger?
\end{quote}

Masculine.

\begin{cousinweb}
\textbf{\ur{پلگ}} (\emph{plag}) is borrowed from English \emph{plug}.
\end{cousinweb}

\begin{grammarlens}[title={What you've built}]
One more word for describing what you lived and felt.
\end{grammarlens}

\subsection*{Your turn}
\begin{itemize}
\raggedright
  \item \emph{Say it:} \emph{plag}
  \item \emph{Say it:} \emph{plag}, once more
  \item \emph{Say it:} \emph{is chārjar kā plag kahāṅ hai?} — Where is the plug for this charger?
  \item \emph{Recall:} say the Urdu for to smear (with mud), then the Urdu for a switch, then say \emph{plag} again
\end{itemize}

\begin{tcolorbox}[breakable,colback=teal!4,colframe=teal!35!black,title={Before you move on}]
What does \ur{پلگ} mean? (\textquotedblleft{}A plug (on a cable)\textquotedblright{}.) Say it once more.
\end{tcolorbox}
\section[\texorpdfstring{\ur{کرنٹ} (karanṭ)}{کرنٹ (karanṭ)}]{\ur{کرنٹ} (karanṭ) — an electric current}

\label{lesson:UR-C504-karant}



\begin{quote}
Before the new one: say the Urdu for a switch, then the Urdu for a plug (on a cable).
\end{quote}

\subsection*{You'll want to know: \ur{کرنٹ}}
\textbf{\ur{کرنٹ}} — \emph{karanṭ} — \textquotedblleft{}an electric current\textquotedblright{}.

\begin{quote}
\ur{اس} \ur{تار} \ur{کو} \ur{نہ} \ur{چھوئیے،} \ur{اس} \ur{میں} \ur{کرنٹ} \ur{ہے۔}
\end{quote}

\begin{quote}
\emph{is tār ko na chhūiye, is meṅ karanṭ hai.}
\end{quote}

\begin{quote}
Please don't touch this wire: there is current in it.
\end{quote}

Masculine.

\begin{cousinweb}
\textbf{\ur{کرنٹ}} (\emph{karanṭ}) is borrowed from English \emph{current}.
\end{cousinweb}

\begin{grammarlens}[title={What you've built}]
One more word for describing what you lived and felt.
\end{grammarlens}

\subsection*{Your turn}
\begin{itemize}
\raggedright
  \item \emph{Say it:} \emph{karanṭ}
  \item \emph{Say it:} \emph{karanṭ}, once more
  \item \emph{Say it:} \emph{is tār ko na chhūiye, is meṅ karanṭ hai.} — Please don't touch this wire: there is current in it.
  \item \emph{Recall:} say the Urdu for a switch, then the Urdu for a plug (on a cable), then say \emph{karanṭ} again
\end{itemize}

\begin{tcolorbox}[breakable,colback=teal!4,colframe=teal!35!black,title={Before you move on}]
What does \ur{کرنٹ} mean? (\textquotedblleft{}An electric current\textquotedblright{}.) Say it once more.
\end{tcolorbox}
\section[\texorpdfstring{\ur{کھونٹی} (khūṅṭī)}{کھونٹی (khūṅṭī)}]{\ur{کھونٹی} (khūṅṭī) — a peg (on a wall)}

\label{lesson:UR-C504-khunti}



\begin{quote}
Before the new one: say the Urdu for a plug (on a cable), then the Urdu for electric current.
\end{quote}

\subsection*{You'll want to know: \ur{کھونٹی}}
\textbf{\ur{کھونٹی}} — \emph{khūṅṭī} — \textquotedblleft{}a peg (on a wall)\textquotedblright{}.

\begin{quote}
\ur{اس} \ur{کا} \ur{تھیلا} \ur{کھونٹی} \ur{پر} \ur{ہے۔}
\end{quote}

\begin{quote}
\emph{us kā thailā khūṅṭī par hai.}
\end{quote}

\begin{quote}
Her bag is hanging on the peg.
\end{quote}

Feminine.

\begin{cousinweb}
\textbf{\ur{کھونٹی}} (\emph{khūṅṭī}) is inherited: CDIAL lists the Hindi cousin \emph{khū̃ṭī}, \textquotedblleft{}a small peg\textquotedblright{}, under \emph{*khuṭṭa}, a form it reconstructs (CDIAL 3893).
\end{cousinweb}

\begin{grammarlens}[title={What you've built}]
One more word for describing what you lived and felt.
\end{grammarlens}

\subsection*{Your turn}
\begin{itemize}
\raggedright
  \item \emph{Say it:} \emph{khūṅṭī}
  \item \emph{Say it:} \emph{khūṅṭī}, once more
  \item \emph{Say it:} \emph{us kā thailā khūṅṭī par hai.} — Her bag is hanging on the peg.
  \item \emph{Recall:} say the Urdu for a plug (on a cable), then the Urdu for electric current, then say \emph{khūṅṭī} again
\end{itemize}

\begin{tcolorbox}[breakable,colback=teal!4,colframe=teal!35!black,title={Before you move on}]
What does \ur{کھونٹی} mean? (\textquotedblleft{}A peg (on a wall)\textquotedblright{}.) Say it once more.
\end{tcolorbox}
\section[\texorpdfstring{\ur{ساکٹ} (sākeṭ)}{ساکٹ (sākeṭ)}]{\ur{ساکٹ} (sākeṭ) — a socket (for a plug)}

\label{lesson:UR-C504-saket}



\begin{quote}
Before the new one: say the Urdu for electric current, then the Urdu for a peg (on a wall).
\end{quote}

\subsection*{You'll want to know: \ur{ساکٹ}}
\textbf{\ur{ساکٹ}} — \emph{sākeṭ} — \textquotedblleft{}a socket (for a plug)\textquotedblright{}.

\begin{quote}
\ur{اس} \ur{کمرے} \ur{میں} \ur{دو} \ur{ساکٹ} \ur{ہیں۔}
\end{quote}

\begin{quote}
\emph{is kamre meṅ do sākeṭ haiṅ.}
\end{quote}

\begin{quote}
There are two sockets in this room.
\end{quote}

Masculine.

\begin{cousinweb}
\textbf{\ur{ساکٹ}} (\emph{sākeṭ}) is borrowed from English \emph{socket}.
\end{cousinweb}

\begin{grammarlens}[title={What you've built}]
One more word for describing what you lived and felt.
\end{grammarlens}

\subsection*{Your turn}
\begin{itemize}
\raggedright
  \item \emph{Say it:} \emph{sākeṭ}
  \item \emph{Say it:} \emph{sākeṭ}, once more
  \item \emph{Say it:} \emph{is kamre meṅ do sākeṭ haiṅ.} — There are two sockets in this room.
  \item \emph{Recall:} say the Urdu for electric current, then the Urdu for a peg (on a wall), then say \emph{sākeṭ} again
\end{itemize}

\begin{tcolorbox}[breakable,colback=teal!4,colframe=teal!35!black,title={Before you move on}]
What does \ur{ساکٹ} mean? (\textquotedblleft{}A socket (for a plug)\textquotedblright{}.) Say it once more.
\end{tcolorbox}
